Scaling quantum processors requires increasing accessible quantum resources whilst retaining sufficient control over them. 
%
The trapped-ion platform shows great promise for realising scaled devices. However, the use of common trapping and control fields can introduce structured multi-body error channels that lead to unintended correlations. 
%
This thesis explores the impact of correlated errors using a recently commissioned experimental platform, including a single-ion addressing system for selective control of the electronic, and motional states on up to 10 ions.

Methods for characterising, simplifying, and suppressing errors are crucial to retain control of scaled systems. 
%
Complete error characterisation is prohibitively expensive with increasing register size. Therefore, aggregated metrics are required to capture relevant noise processes.

To quantify correlation-inducing errors on scaled systems, we first develop \emph{k-copy} randomised benchmarking. This method utilises symmetries enforced by structured noisy fields to discern correlated and local error components acting on a multi-qubit devices.
%
We use this method to resolve spatially correlated errors caused by a global-addressing control field acting on an eight-ion chain.
%
Such correlated errors are particularly impactful to error correction and mitigation techniques that assume local, depolarising noise.

In a second experiment, we explore the impact of correlated errors on virtual distillation, an error mitigation protocol using multiple identical noisy copies to estimate error-suppressed expectation values. A suppression of $\times2.8(7)$ is measured when subject to dominant local noise, whilst under dominant common-mode correlating errors, the suppression is reduced to $\times1.20(9)$. We further show that purely common-mode errors are entirely unaffected by virtual distillation.
%
These results highlight the importance of resolving error structure, not simply overall error rates on the constituent subsystems. 

Finally, we explore a complementary route to scaling the available quantum resources: using the motion of the ions for information processing. 
%
Building on previous demonstrations of hybrid qubit-qumode control in trapped-ions, we demonstrate an elementary method to generate spin-conditioned bosonic interactions, now using the single-ion addressing system. This enables scaling to multiple ions, and access to many qumodes.
%
The work presented in this thesis develops an experimental platform for studying larger qubit-qumode registers, and methods to characterise and mitigate emergent correlated errors.